

Algebraic topology has its roots in constructing, computing and understanding the deeper structure of three main invariants classicaly associated to a topological space $X$: homotopy groups $\pi_* X$, homology groups $H_*X$, and cohomology groups $H^*X$. Elements of the $n$-th homotopy group $\pi_n X$ are classified by certain homotopy classes of maps $[S^n, X]$ from the $n$-dimensional sphere $S^n$ into your space of choice $X$, and often tell you more about a space than (co)homology groups. The tradeoff is that homotopy groups are in general much more difficult to compute, even providing a full description of the homotopy groups of spheres is an elusive and exceptionally difficult problem.

As a result, it can be useful to know even an approximation of the homotopy groups of a space. The \textit{stable homotopy groups} $\pi_*^sX$ of a space $X$ are one such approximation. It follows from the Freudenthal suspension theorem that if $X$ is a $k$-connected space\footnote{This means that the groups $\pi_n X$ for $n\leq k$ are trivial} that $\pi_n X\cong \pi_n^sX$ for $n\leq 2k$; i.e., that the stable homotopy groups agree with the homotopy groups in a ``stable range''. For various reasons (which we will describe in more detail shortly), these stable homotopy groups are often more computable than ordinary homotopy groups, which allows one to make meaningful progress by working with the easier-to-compute stable homotopy groups instead.

But how are the stable homotopy groups constructed? The reader familiar with homology theory will recall that there is always an isomorphism \begin{equation} \label{eq:homology-iso} H^*X\cong H^{*+1}(\Sigma X)\end{equation}  where $\Sigma X$ denotes the \textit{suspension} of a space $X$ (which may be modelled as the smash product in spaces of the $1$-sphere and $X$, $\Sigma X=S^1\wedge X$). Freudenthal's suspension theorem tell us that, unlike homology groups, the homotopy groups of a space only obey this isomorphism inside the stable range, i.e., $\pi_nX\cong \pi_{n+1}(\Sigma X)$ for $n\leq 2k$ when $X$ is $k$-connected. However, in spaces the smash product $\wedge$ admits a right adjoint $Y\mapsto \Omega Y=\Map(S^1,Y)$ \footnote{The latter here is topologized by the compact-open topology, though $\Omega Y$ can more easily be thought of as the \textit{space} of loops in $Y$ based at $*\in Y$.}. Using the identification $\pi_{n+1}(\Sigma X)\cong [S^{n+1}, S^1 \wedge X]$, the isomorphism above is obtained by applying the $n$-th homotopy group functor to the \textit{point-set level} map $\eta_X \colon X\to \Omega \Sigma X$. Iterating this map builds the \textit{stabilization} $QX$ of a topological space $X$ as the following colimit whose maps are given by $\Omega^n \eta_{\Sigma^n X}$ \[ QX:= \U \Sp X=\colim( X\to \Omega \Sigma X\to \Omega^2 \Sigma^2 X\to \Omega^3 \Sigma^3 X \to \cdots).\]

The stable homotopy groups are obtained then as the ordinary homotopy groups of $QX$, i.e., \[\pi_*^sX:= \pi_*(QX).\] What's more is that, in the language of Goodwillie's calculus of homotopy functors, the functor $Q$ may be thought of as a \textit{first order}---or \textit{linear}---approximation to the identity functor on based topological spaces. This strikes the question: What does linearity mean in this context?

Inspiration may be taken again from homology. Again, the familiar reader may recall that homology groups satisfy \textit{excision} which gives rise to a strong computational tool, the Mayer-Vietoris sequence. An equivalent statement is that any pushout square \[ \xymatrix{A \ar[r] \ar[d] & C \ar[d] \\ B \ar[r] & D }\] of CW complexes where all of the maps are inclusions, be taken to a pullback square of groups \[ \xymatrix{H_* A \ar[r] \ar[d] & H_* C \ar[d] \\ H_* B \ar[r] & H_* D. }\] By setting $C$ and $D$ to be the cone on $A$, and $D=\Sigma A$, the isomorphism \eqref{eq:homology-iso} follows from a simple calculation. 

The functor $Q$ satisfies a related homotopical condition: that any \textit{homotopy} pushout square of spaces \[ \xymatrix{X \ar[r] \ar[d] & W \ar[d] \\ Y \ar[r] & Z }\] be taken to a \textit{homotopy} pullback square \[ \xymatrix{QX \ar[r] \ar[d] & QW \ar[d] \\ QY \ar[r] & QZ. }\] Such functors are called \textit{excisive}, or more specifically \textit{$1$-excisive}. In \cite{Goo-Calc3}, Goodwilllie shows that any functor $F\colon \Top\to \Top$ which preserves weak homotopy equivalences admits a universal such $1$-excisive approximation, $P_1F$, by means of a comparison map $F\to P_1F$ (which, under suitable niceness conditions on $F$ should induce an isomorphism in a ``stable range'' of homotopy groups when evaluated on suitably connected spaces). Indeed, the universal $1$-excisive approximation to the identity functor on $\Top$ is the stable homotopy functor $Q$. 

Let $\Spt$ denote (a model for) the category of \textit{spectra} such that $\Spt$ is symmetric monoidal with respect to the smash product $\wedge$, for which the sphere spectrum $S$ is the unit. The term \textit{linear} can be justified by observing that any $1$-excisive functor $\Spt\to \Spt$ is of the form $X\mapsto E\wedge X$ for fixed spectrum $E$ (at least when restricted to \textit{finite} spectra $X$). Moreover, $1$-excisive functors $F\colon \Top\to \Top$ such that $F(\pt)\simeq \pt$ are essentially those which factor through the the usual \textit{stabilization adjunction} $(\Sp, \U)$ between $\Top$ and $\Spt$, \[ \xymatrix{\Top \ar@<1ex>[r]^-{\Sp} & \Spt \ar@<.5ex>[l]^-{\U}}, \] as $F\simeq \U \widetilde{F} \Sp$ such that $\widetilde{F}$ is a $1$-excisive functor on spectra. 

In fact, Goodwillie shows more. In \cite{Goo-Calc3}, he constructs a \textit{Taylor tower} of $n$-excisive or ``polynomial of degree at most $n$'' functors, $P_n F$, along with and natural transformations $P_{n+1} F\to P_{n} F$ (for $n\geq 0$), of the form
\[
   F\to \cdots \to P_n F\to P_{n-1} F\to \cdots \to P_1 F\to P_0 F. %give number to this
\]
These functors $P_n F$ play a role analogous to the Taylor polynomials $p_nf$ associated to an infinitely differentiable function $f\colon \mathbf{R}\to\mathbf{R}$ in ordinary calculus and analysis. The formula \[ p_n f(x)-p_{n-1} f(x)=\dfrac{f^{(n)}(0)}{n!} x^n\] has a homotopical analog in that the ``difference'' (in this case, \textit{homotopy fiber}) between $P_n F$ and $P_{n-1}F$ is measured by a certain spectrum with action by the $n$-th symmetric group $\Sigma_n$ denoted $\partial_n F$ (the \textit{$n$-th derivative} of $F$) as an equivalence of the form \begin{equation*} \label{eq:homog-layer} D_nF(X) :=\hofib (P_n F(X) \to P_{n-1} F(X)) \simeq \U (\partial_n F \wedge \Sp X^{\wedge n})_{h \Sigma_n}. \end{equation*} The subscript $h\Sigma_n$ above denotes that homotopy orbits have been taken with respect to the $\Sigma_n$ action on $\partial_n F$ and on $\Sp X^{\wedge n}$ (via permutation of factors). 

Under suitable niceness conditions on $F$ and the input space $X$, the Taylor tower may be used to tell us something about the homotopy type of $F(X)$. For instance, the Taylor tower of the identity on $\Top$ will converge to the homotopy type of its input when evaluated on simply-connected spaces. That is, for $1$-connected $X\in \Top$, the natural map \[ X\to \lim_{n\to \infty} P_n \Id_{\Top} (X)\] is a weak homotopy equivalence. As such, the Taylor tower of the identity provides a useful canonical resolution of a space which begins with the stabilization $QX$. Moreover, this resolution is ``controlled'' by the symmetric sequence of derivatives $\partial_* \Id_{\Top}=\{\partial_n \Id_{\Top} \}_{n\geq 1}$

Much work has gone into understanding Taylor towers and Goodwillie derivatives in recent decades, both for functors of spaces and spectra and also in more general contexts (see \cite{Per-thesis}, \cite[\textsection 6]{Lur}). Johnson \cite{Joh} and Arone-Mahowald \cite{Ar-Mah} give a description of the layers $D_n \Id_{\Top}$ in terms of Spanier-Whitehead duals to the $n$-th \textit{Partition poset complex}---essentially, a pointed space built from the combinatorics of partitions of the set $\{1,\dots,n\}.$ Using this description, Ching \cite{Ching-thesis} has shown that the symmetric sequence of Goodwillie derivatives $\partial_* \Id_{\Top}=\{\partial_n \Id_{\Top} \}_{n\geq 1}$ can be given a natural \textit{operad} structure. This operad plays a central rule in describing a \textit{chain rule} for derivatives as shown by Arone-Ching \cite{AC-1}: that is, an equivalence of symmetric sequences \[ \partial_*F \circ_{\partial_*\Id_{\Top}} \partial_*G\simeq \partial_*(FG)\] for functors $F,G\colon \Top\to \Top$ (here $\circ$ denotes the \textit{composition product}, see \ref{comp-prod-def-orb}). 

This operad of Goodwillie derivatives indeed plays a central role in describing the homotopy theory of the category of based spaces: In particular, Heuts \cite{Heuts-vn} (see also \cite{BR}) has further shown that certain ``chromatic localizations'' $M_n^f$ of $\Top$ can be characterized by the category of algebras over $\partial_*\Id_{\Top}$ in $T(n)$-local spectra. As such, it is often anticipated that similar techniques may be used to better understand other categories in which one can \textit{do} functor calculus. A widely recognized slogan of functor calculus that the Goodwillie derivatives of the identity on a suitable model category $\mathsf{C}$ should come equipped with a canonical operad structure which tells us something about $\mathsf{C}$. One such example, and the focus of our first main theorem, is when $\mathsf{C}$ is a category of algebras over a (reduced) operad $\Cal{O}$ in spectra. 


First, some background. Operads \cite{May-geom}, \cite{BV} are combinatorial tools for describing spaces (or spectra) $X$ which admit a pairing $X\wedge X\to X$ and unit map $S\to X$ that may be only associative or commutative up to coherent homotopies (such as the structure found on an $n$-fold loop space $\Omega^n Y$ \cite{May-geom}) and have played an increasingly common and powerful role in homotopy theory. Common examples of interest are associative, commutative, or $E_n$-ring spectra (for $1\leq n\leq \infty)$---the latter of which interpret between \textit{highly homotopy associative} ($E_1$ or $A_{\infty}$) and \textit{highly homotopy commutative} ($E_{\infty}$) ring spectra and enjoy a rich structure on the (co)homology theories they represent. For an operad $\Cal{O}$ in spectra, we write $\Alg{O}$ for the category of algebras over $\Cal{O}$. An object $X\in \Alg{O}$ is a spectrum $X$ together with associative and unital action maps of the form \[ \Cal{O}[n] \wedge_{\Sigma_n} X^{\wedge n}\to X \quad \quad (n\geq 0)\] We will require that our operads be \textit{reduced} in that the $0$-ary operations described by $\Cal{O}$ are trivial (i.e., $\Cal{O}[0]=\pt$). Such condition guarantees that our $\Cal{O}$-algebras are \textit{nonunital}; a condition that can naturally arises when working with augmented rings. 

Pereira \cite{Per-thesis} shows that the constructions and main results of \cite{Goo-Calc3} readily extend to functors between categories of algebras over operads in spectra. Many other authors are responsible for early work which suggests this should be possible including Kuhn \cite{KuhFC}, McCarthy-Minasian \cite{MM}, Basterra-Mandell \cite{Bast-Man}, and Harper-Hess \cite{Har-Hess}. Our work relies heavily on the constructions of Kuhn-Pereira \cite{KuhPer} for describing the Taylor tower of certain functors on $\Alg{O}$. 


\section{Main theorems}  Our first main result, Theorem \ref{main}(a) as follows, is that the Goodwillie derivatives of the identity in the category of algebras over a reduced operad $\Cal{O}$ in spectra can be given a naturally occurring ``highly homotopy coherent'' operad structure. 

\begin{theorem} \label{main} Let $\Cal{O}$ be an operad in spectra such that $\Cal{O}[n]$ is $(-1)$-connected for $n\geq 1$ and $\Cal{O}[0]=\pt$. Then,
    \begin{enumerate}[label=(\alph*)]
        \item The derivatives of the identity in $\Alg{O}$ can be equipped with a natural highly homotopy coherent operad structure.
        \item Moreover, with respect to this structure, $\partial_* \Id_{\Alg{O}}$ is equivalent to $\Cal{O}$ as highly homotopy coherent operads. 
    \end{enumerate}
\end{theorem}

It has been known for some time that $\Cal{O}[n]$ is a model for $\partial_n \Id_{\Alg{O}}$. Part (b) of Theorem \ref{main} in particular answers a long standing conjecture which appears in \cite{AC-1} (and also answered in the context of $\infty$-operads by Ching in \cite{Ching-day-conv}), a main difficulty of which is describing an intrinsic operad structure on the derivatives of the identity which may be compared with that of the operad $\Cal{O}$.  

The proofs of parts (a) and (b) to Theorem \ref{main} may be found in Sections \ref{sec:1a} and \ref{O-and-dern-equiv}, respectively. Our technique is to avoid working with the identity directly by replacing it with the Bousfield-Kan cosimplicial resolution provided by the stabilization adjunction $(Q,U)$ for $\Cal{O}$-algebras. The strong cartesianness estimates of Blomquist \cite{Blom} (see also Ching-Harper \cite{Ching-Harper-BM}) allow us to then express $\partial_* \Id_{\Alg{O}}$ as the homotopy limit of the cosimplicial diagram (showing only coface maps) \begin{equation} \label{eq:der-model-stab} \partial_*(QU)^{\bullet+1}=  \left(\xymatrix{
    \partial_*(UQ) \ar@<-.5ex>[r] \ar@<.5ex>[r] &
    \partial_*(UQ)^2 \ar@<-1ex>[r] \ar[r] \ar@<1ex>[r] &
    \partial_*(UQ)^3  \cdots 
}\right)\end{equation}
whose terms $\partial_*(QU)^{k+1}$ may be readily computed by an $\Cal{O}$-algebra analogue of the Snaith splitting. We thus obtain a natural cosimplicial resolution $C(\Cal{O})$ of the derivatives of the identity such that $\partial_* \Id_{\Alg{O}}\simeq \holim_{\Delta} C(\Cal{O})$ which furthermore may be identified as the $\TQ$ resolution of $\Cal{O}$ as a left $\Cal{O}$-module. Our approach is influenced by the work of Arone-Kankaanrinta \cite{AK} wherein they use the cosimplicial resolution offered by the stabilization adjunction between spaces and spectra to analyze the derivatives of the identity in spaces via the classic Snaith splitting. 

We induce a highly homotopy coherent operad structure (i.e., $A_{\infty}$-operad) on $\partial_* \Id_{\Alg{O}}$ by constructing a pairing of the resolution $C(\Cal{O})$ with respect to the box product $\square$ for cosimplicial objects (see Batanin \cite{Bat-box}). Thus, we extend to the monoidal category of symmetric sequences a technique utilized in McClure-Smith \cite{MS-box}: specifically, that if $X$ is a $\square$-monoid in cosimplicial spaces or spectra then $\Tot(X)$ is an $A_{\infty}$-monoid (with respect to the closed, symmetric monoidal product for spaces or spectra).

There are some subtleties that arise in that (i) the box product is not as well-behaved when working with the composition product $\circ$ of symmetric sequences, and (ii) the extra structure encoded by $\circ$ leads us to work with $\mathbf{N}$-colored operads to express $A_{\infty}$-monoids with respect to composition product.  As such, one of the main developments of this thesis is that of $\mathbf{N}$-colored operads \textit{with levels} (i.e., $\Nl$-operads) as useful bookkeeping tools designed to algebraically encode operads (i.e., strict composition product monoids) and ``fattened-up'' operads as their algebras. Within this framework of $\Nl$-operads we can also describe algebras over an $A_{\infty}$-operad.

\begin{remark} \label{sec:remark}
In the statement of Theorem \ref{main} the phrase ``naturally occuring'' means that we refrain from endowing $\partial_* \Id_{\Alg{O}}$ with the operad structure from $\Cal{O}$ directly. Rather, we produce a method for intrinsically describing operadic structure possessed by the derivatives of the identity that should carry over to other model categories suitable for functor calculus. In particular, the constructions of such an operad structure on the derivatives of the identity should:

 \begin{enumerate}[label=(\roman*)] 
\item Recover the ($A_{\infty}$-) operad structure endowed on $\partial_* \Id_{\mathsf{Top}_*}$ described by Ching in \cite{Ching-thesis}.

\item Endow the derivatives of an suitably nice homotopy functor  \footnote{For instance, we should expect this to be possible for functors which are finitary and simplicial} $F\colon \Alg{O}\to \Alg{O'}$ with a natural $(\partial_* \Id_{\Alg{O'}}, \partial_* \Id_{\Alg{O}})$-bimodule structure (which is in turn equivalent to an $(\Cal{O}',\Cal{O})$-bimodule) suitable for describing a chain rule (as in Arone-Ching \cite{AC-1}) with a view toward running the machinery of \cite{AC-coalg}.

\item Be fundamental enough to describe an operad structure on $\partial_* \Id_{\mathsf{C}}$ and chain rule for a suitable model category $\mathsf{C}$ (e.g.,  one in which one can do functor calculus).
\end{enumerate}

In fact, item (i) is the subject of our second main result, Theorem \ref{thm:main-1} below. Items (ii) and (iii) are the subject of current work, and we defer a discussion to our progress along with some conjectural remarks to Section \ref{sec:conj}. \end{remark}

 \begin{theorem} \label{thm:main-1}
The symmetric sequence $\partial_*\Id_{\Top}$ is a highly homotopy coherent operad.
\end{theorem}

One of our main interests in operads is in what structure is described on their algebras. Ching further shows that the homology of $\partial_*\Id_{\Top}$ gives a ``Lie operad'' \cite{Ching-thesis}, which is \textit{Koszul dual} to the commutative cooperad \cite{Ginz-Kap} (see also \cite{Get-Jon}). The operad $\partial_*\Id_{\Top}$ of spectra is similarly Koszul dual to the commutative cooperad in spectra (as in \cite{FG}) and so is often referred to as the \textit{spectral Lie operad}.

In addition, we prove the following theorem (Theorem \ref{thm:main-2}). As the commutative cooperad in $\Spt$ is also an operad, Theorem \ref{thm:main-2} may be thought of as an amplified version of the classic result from commutative algebra which states that the primitives of a Hopf algebra naturally form a Lie algebra (see, e.g. \cite{Abe}). 

\begin{theorem} \label{thm:main-2}
    The derived primitives (Definition \ref{def:derived-prim}) of a commutative coalgebra \footnote{For us, coalgebra means \textit{coalgebra with divided powers} (see, e.g. \cite{FG}).} in spectra admit a natural action by $\partial_*\Id_{\Top}$. 
\end{theorem}

The above theorem is also not necessarily new. For a finite commutative coalgebra $Y$, the derived primitives $\Prim(Y)$ is equivalent to $\TQ(Y^{\vee})^{\vee}$ (see Example \ref{rem:SW-TQ})---where here $Y^{\vee}=\Map(Y,S)$ is the Spanier-Whitehead dual---and the latter is known to admit an action by $\partial_* \Id_{\Top}$ (e.g. as in \cite{Ching-thesis}, \cite{BR}, \cite{Heuts-vn}).  However, our constructions provide an easy to use, algebraic alternative to the pioneering operad structure on $\partial_*\Id_{\Top}$ found in \cite{Ching-thesis}. In particular, we expect our approach to allow us to show that the Goodwillie derivatives of the identity in a suitable model category $\mathsf{C}$ come equipped with a canonical operad structure, induced in a similar way. 


\subsection{Outline of the argument}
Our main tool is to utilize the Bousfield-Kan cosimplicial resolution with respect to the stabilization. For the context of $\Cal{O}$-algebras in Theorem \ref{main}, this relies on a result of Basterra \cite{Bast}, that the stabilization of an $\Cal{O}$-algebra $X$ is naturally equivalent to its \textit{topological Quillen homology} spectrum $\TQ(X)$ (see Section \ref{sec:stab}). Topological Quillen homology may be thought of as the derived indecomposables quotient for structured ring spectra described as algebras over a (reduced) operad, and arises as the left derived functor of the left adjoint in \[ \xymatrix{
    \Alg{O} \ar@<.75ex>[r]^-{Q} & \textsf{Alg}_{J}\simeq \mathsf{Mod}_{\Cal{O}[1]}. \ar@<.5ex>[l]^-{U}
    }\] 
Here, $J$ denotes a suitable replacement of $\tau_1\Cal{O}$, the truncation of $\Cal{O}$ above level 1 (see Section \ref{sec:stab}), $U$ denotes the forgetful functor along the map of operads $\Cal{O}\to J$ (which also provides trivial operations from $\Cal{O}[n]$ on a $J$-algebra for $n\geq 2$), and $Q$ is the $\Cal{O}$-algebra analog of the indecomposables quotient $R/R^2$ for $R$ a nonunital commutative ring in classical algebra. 

Using the strong connectivity estimates offered by Blomquist's higher stabilization theorems \cite[\textsection 7]{Blom}, we first show that $\partial_* \Id_{\Alg{O}}$ is equivalent to $\holim_{\Delta} \partial_*(UQ)^{\bullet+1}$ (see (\ref{eq:der-model-stab})). Similarly to Arone-Kankaanrinta \cite{AK}, in which they compute the $n$-excisive approximations (resp. $n$-th derivatives) of the identity functor on $\Top$ in terms of the $n$-excisive approximations (resp. $n$-th derivatives) of iterates of stabilization $\U \Sp$ by means of the Snaith splitting, we then analyze the terms $\partial_*(UQ)^{k+1}$ via an analog of the Snaith splitting in $\Alg{O}$.

Essentially a statement about the Taylor tower of the associated comonad $QU$, the Snaith splitting in $\Alg{O}$ permits equivalences of symmetric sequences\[\partial_*(QU)\simeq  |\Barr(J, \Cal{O}, J)|\simeq J\circ^{\mathsf{h}}_{\Cal{O}} J =:B(\Cal{O})\] as $(J,J)$-bimodules (here, $\circ^{\mathsf{h}}$ denotes the derived composition product). By iterated applications of the splitting, we may compute \[\partial_*(UQ)^{k+1}\simeq \underbrace{B(\Cal{O}) \circ_J \cdots \circ_J B(\Cal{O})}_{k}\simeq \underbrace{J\circ_{\Cal{O}}\cdots  \circ_{\Cal{O}} J}_{k+1}= C(\Cal{O})^k\] and moreover that $\partial_*(UQ)^{\bullet+1}\simeq C(\Cal{O})$ as cosimplicial symmetric sequences. Here, $C(\Cal{O})$ is given by \[\xymatrix{
    J \ar@<.5ex>[r] \ar@<-.5ex>[r] &
    J\circ_{\Cal{O}}J \ar@<-1ex>[r] \ar[r] \ar@<1ex>[r] &
    J\circ_{\Cal{O}}J \circ_{\Cal{O}}J \ar@<-1.5ex>[r] \ar@<-.5ex>[r] \ar@<.5ex>[r] \ar@<1.5 ex>[r] &
    J\circ_{\Cal{O}}J \circ_{\Cal{O}}J\circ_{\Cal{O}} J \cdots
}\]
with coface map $d^i$ induced by inserting $\Cal{O}\to J$ at the $i$-th position (see Remark \ref{prop:der_n-model-equivs} along with (\ref{TQ-res-1})). 

Note, $B(\Cal{O})$ is (at least up to homotopy) a cooperad with a coaugmentation map $J\to B(\Cal{O})$, and our $C(\Cal{O})$ is essentially a rigid cosimplicial model for the cobar construction on $B(\Cal{O})$. In particular, this allows us to bypass referencing any particular model for the comultiplication on $B(\Cal{O})$ (e.g.,  that of Ching \cite{Ching-thesis}, see also Section \ref{B(O)}).  

We construct a pairing $m\colon C(\Cal{O})\square C(\Cal{O})\to C(\Cal{O})$ with respect to the box product (Definition \ref{def:box-product}) of cosimplicial symmetric sequences via compatible maps of the form (induced by the operad structure maps $J\circ J\to J$) \[ m_{p,q}\colon \underbrace{J\circ_\Cal{O} \cdots \circ_{\Cal{O}}(J}_{p+1}\circ \underbrace{J)\circ_\Cal{O} \cdots \circ_{\Cal{O}}J}_{q+1} \to \underbrace{J\circ_\Cal{O}\cdots \circ_{\Cal{O}} J\circ_{\Cal{O}} \cdots \circ_{\Cal{O}}J}_{p+q+1}\] along with a unit map $u\colon \underline{I}\to C(\Cal{O})$, where $\underline{I}$ denotes the constant cosimplicial symmetric sequence on $I$. Our argument is then to induce an $A_{\infty}$-monoidal pairing on $\partial_* \Id_{\Alg{O}}$ --- modeled as $\Tot C(\Cal{O})$ --- via $m$ and $u$ (compare with McClure-Smith \cite{MS-box}).

One difficulty which arises is that the composition product of symmetric sequences is not as well-behaved of a product as, say, cartesian product of spaces or smash product of spectra. Thus, we \textit{do not} obtain $m$ as a strictly monoidal pairing on the level of cosimplicial diagrams. In resolving this issue we introduce a specialized category of $\mathbf{N}$-colored operads \textit{with levels} (i.e., $\Nl$-operads) designed specifically to overcome these technical subtleties of the composition product. As a result, a large portion of this document is dedicated to carefully developing the framework of $\Nl$-operads and their algebras. 

With these details in tow it is then possible to produce an $A_{\infty}$-operad structure on $\partial_* \Id_{\Alg{O}}$. Let $\Tot$ denote restricted totalization $\mathrm{Tot}^{\mathsf{res}}$ (see Section \ref{sec:Use-of-tot}), we then obtain an $A_{\infty}$-monoidal pairing \[\Tot C(\Cal{O})\circ \Tot C(\Cal{O})\to \Tot C(\Cal{O})\] described as an algebra over a certain $\Nl$-operad which is a naturally ``fattened-up'' replacement of the $\Nl$-operad whose algebras are strict operads (see Definition \ref{def-of-Oper} along with Propositions \ref{prop:Oper-is-Nl} and \ref{oper-oper}). Moreover, the coaugmentation $\Cal{O}\to C(\Cal{O})$ provides a comparison between $\Cal{O}$ and $\partial_* \Id_{\Alg{O}}$ which we show yields an equivalence of $A_{\infty}$-operads, thus resolving the aforementioned conjecture. 

Our method for proving Theorems \ref{thm:main-1} and \ref{thm:main-2} is similar. We make use of the models for $\partial_n \Id_{\Top}$  as the Spanier-Whitehead dual of the $n$-th partition poset complex $\Par(n)$ (Definition \ref{def:par-poset}), \[ \partial_n \Id_{\Top} \simeq \Map(\Par(n),S), \quad \quad (n\geq 1).\]

We then show that $\Map(\Par(n),S)$ can be modeled as the totalization of the cobar resolution $C(\uS)[n]$ on the commutative cooperad of spectra. Essentially, $C(\uS)[n]$ has $p$-simplices given by the $p$-fold composition $\uS^{\circ p}[n]$, where $\uS$ is the reduced symmetric sequence with $\uS[n]=S$ (with trivial $\Sigma_n$ action) for $n\geq 1$. We show that the cosimplicial symmetric sequence $C(\uS)$ admits a natural $\boxcirc$-monoid structure (see \cite[Definition 5.6]{Clark-1}) induced, essentially, by the canonical isomorphisms $\uS^{\circ p} \circ \uS^{\circ q} \cong \uS^{\circ p+q}$. This $\boxcirc$-monoid structure then induces the desired ``highly homotopy coherent'' operad structure on $\partial_* \Id_{\Top}$ upon passing to totalization. 

We similarly show that the derived primitives may be calculated as the totalization of a certain cosimplicial spectrum $C(Y)$. In particular, $C(Y)$ is precisely the dual notion of derived indecomposables which underlies the construction of topological Quillen homology for $\Cal{O}$-algebras. We show that $C(Y)$ is a $\boxcirc$-module over the cobar resolution $C(\uS)$ and that this module structure induces an action of $\partial_*\Id_{\Top}$ on the derived primitives upon passing to totalization. 


\subsection{Comparison to the operad structure constructed by Ching} We expect the operad structure on $\partial_* \Id_{\Top}$ described in this document is equivalent to that constructed by Ching in \cite{Ching-thesis}. 

In particular, our main technical lemma (Lemma \ref{lem:tech}) is reminiscent of the ``tree ungrafting'' arguments found in \cite{Ching-thesis}, and our notion of $\boxcirc$-monoid is similar to Ching's notion of ``pre-cooperad'' in \cite{Ching-barcobar}. However, despite these similarities, the author is not aware of an explicit comparison between the two constructions, nor to the third description offered in \cite{Ching-day-conv}.